\documentclass[11pt,a4paper]{article}
\usepackage[utf8]{inputenc}
\usepackage[T1]{fontenc}
% \usepackage[document]{ragged2e}
% \usepackage[a4paper, total={5in, 8in}]{geometry}
\usepackage[english]{babel}
\usepackage{amsmath, amsthm, amssymb, amsfonts}
\usepackage{enumerate}
\usepackage{graphicx}
\usepackage{fancyhdr}
\usepackage{booktabs}
\usepackage{float}
\usepackage[colorlinks=true,breaklinks=false]{hyperref}
\usepackage{listings}
\usepackage[norelsize,ruled,vlined]{algorithm2e}
\usepackage{csquotes}
\usepackage{url}
\usepackage{ifthen}
\usepackage{keyval}
\usepackage{etoolbox}
\usepackage{etex}
\usepackage{datetime}
\usepackage{eurosym}
\usepackage{tikz}
\usepackage{wrapfig}
\usepackage{paralist}
\usepackage[backend=bibtex]{biblatex}
\usepackage{xcolor}
\usepackage{caption}
\DeclareCaptionFont{white}{\color{white}}
\DeclareCaptionFormat{listing}{\colorbox[cmyk]{0.43, 0.35, 0.35,0.01}{\parbox{\textwidth}{\hspace{15pt}#1#2#3} } }
\captionsetup[lstlisting]{format=listing, labelfont=white, textfont=white, singlelinecheck=false, margin=0pt, font={bf,footnotesize} }

\hyphenpenalty=9999
% \pretolerance=9999
% \tolerance=9999
% \emergencystretch=0pt
\addbibresource{project.bib}
\newcommand{\scalafmt}{\texttt{scalafmt}}

% Single spacing after dots and colons
\frenchspacing
\hyphenation{scalafmt}
\hyphenation{ClangFormat}
% \hyphenation{DSLVirtualization}
% \hyphenation{ReificationTransformation}

%%%%%%%%%%%%%%%%%%%%%%%%%%%%%%%%%%%%%%%%%%%%%%%%%%%%%%%%%%%%%
%                        Setup
%%%%%%%%%%%%%%%%%%%%%%%%%%%%%%%%%%%%%%%%%%%%%%%%%%%%%%%%%%%%%

\begin{document}
% \raggedright


\definecolor{dkgreen}{rgb}{0,0.6,0}
\definecolor{gray}{rgb}{0.5,0.5,0.5}
\definecolor{mauve}{rgb}{0.58,0,0.82}

\lstdefinestyle{scala}{
  language=Scala,
  aboveskip=3mm,
  basicstyle=\scriptsize\ttfamily,
  belowskip=3mm,
  breakatwhitespace=false,
  showspaces=false,
  showlines=true,
  columns=flexible,
  keepspaces=true,
  commentstyle=\color{dkgreen},
  keywordstyle=\color{blue},
  numbers=left,
  numberstyle=\tiny\color{gray},
  showstringspaces=false,
  stepnumber=1,
}


\definecolor{dkgreen}{rgb}{0,0.6,0}
\definecolor{gray}{rgb}{0.5,0.5,0.5}
\definecolor{mauve}{rgb}{0.58,0,0.82}

\lstset{language=Scala,
  aboveskip=3mm,
  basicstyle={\small\ttfamily},
  belowskip=3mm,
  breakatwhitespace=true,
  breaklines=true,
  columns=flexible,
  commentstyle=\color{dkgreen},
  keywordstyle=\color{blue},
  numbers=left,
  numberstyle=\tiny\color{gray},
  showstringspaces=false,
  stepnumber=1,
  stringstyle=\color{mauve},
  tabsize=2,
}

\usetikzlibrary{calc,trees,positioning,arrows,chains,shapes.geometric,%
    decorations.pathreplacing,decorations.pathmorphing,shapes,%
    matrix,shapes.symbols}
\tikzset{>=stealth',
  punktchain/.style={rectangle,
    rounded corners,
    % fill=black!10,
    draw=black, very thick,
    text width=12.5em,
    minimum height=2.5em,
    text centered,
    on chain},
  every join/.style={->, shorten >=2pt},
}


% Example of title page for the projects carried out within the lasec

% Simply include it in your mastex tex file:
%        % Example of title page for the projects carried out within the lasec

% Simply include it in your mastex tex file:
%        % Example of title page for the projects carried out within the lasec

% Simply include it in your mastex tex file:
%        \input{cover}


% Updated March 2006 (SP)


\newcommand{\logoepfl}[0]{
  \begin{center}
    \includegraphics[width=4cm]{logo_epfl_coul.eps}
  \end{center}
  \vspace{0.3cm}
  \hrule
}
\newcommand{\logolasec}[0]{
  \vspace{1cm}
  \hrule
  \begin{center}
    \includegraphics[width=4.5cm]{logo_lasec_coul.eps}
  \end{center}
}
\newcommand{\project}[1]{
  \begin{center}
    \large{#1}
  \end{center}
  \vspace{1cm}
}
\newcommand{\department}[1]{
  \begin{center}
    \large{#1}
  \end{center}
}
\newcommand{\supervisor}[3]{
  \begin{center}
    \begin{normalsize}{
        \bf #1}\\#2\\#3
    \end{normalsize}
  \end{center}
}
\renewcommand{\author}[1]{
  \begin{center}
    \Large{#1}
  \end{center}
  \vspace{0.5cm}
}
\renewcommand{\title}[1]{
  \vspace{3cm}
  \begin{center}
    \huge{#1}
  \end{center}
  \vspace{1.7cm}
}
\renewcommand{\date}[2]{
  \begin{center}
    \normalsize{#1 #2}
  \end{center}
  \vspace{0.5cm}
}


\thispagestyle{empty}


% begin title page
  \logoepfl

  \title{scalafmt: opinionated \\ code formatter for Scala}

  \author{Ólafur Páll Geirsson}
  \department{School of Computer and Communication Sciences}
  \project{A thesis submitted for the degree of Master of Computer Science at École polytechnique fédérale de Lausanne}

  \date{June}{2016}

  \begin{center}
    \begin{tabular}{cc}
      \begin{tabular}{p{4.0cm}}
        \supervisor{Responsible}{Prof. Martin Odersky}{EPFL / LAMP}
      \end{tabular}&
      \begin{tabular}{p{4.0cm}}
        \supervisor{Supervisor}{Eugene Burmako}{EPFL / LAMP}
      \end{tabular}
    \end{tabular}
  \end{center}

% end title page




% Updated March 2006 (SP)


\newcommand{\logoepfl}[0]{
  \begin{center}
    \includegraphics[width=4cm]{logo_epfl_coul.eps}
  \end{center}
  \vspace{0.3cm}
  \hrule
}
\newcommand{\logolasec}[0]{
  \vspace{1cm}
  \hrule
  \begin{center}
    \includegraphics[width=4.5cm]{logo_lasec_coul.eps}
  \end{center}
}
\newcommand{\project}[1]{
  \begin{center}
    \large{#1}
  \end{center}
  \vspace{1cm}
}
\newcommand{\department}[1]{
  \begin{center}
    \large{#1}
  \end{center}
}
\newcommand{\supervisor}[3]{
  \begin{center}
    \begin{normalsize}{
        \bf #1}\\#2\\#3
    \end{normalsize}
  \end{center}
}
\renewcommand{\author}[1]{
  \begin{center}
    \Large{#1}
  \end{center}
  \vspace{0.5cm}
}
\renewcommand{\title}[1]{
  \vspace{3cm}
  \begin{center}
    \huge{#1}
  \end{center}
  \vspace{1.7cm}
}
\renewcommand{\date}[2]{
  \begin{center}
    \normalsize{#1 #2}
  \end{center}
  \vspace{0.5cm}
}


\thispagestyle{empty}


% begin title page
  \logoepfl

  \title{scalafmt: opinionated \\ code formatter for Scala}

  \author{Ólafur Páll Geirsson}
  \department{School of Computer and Communication Sciences}
  \project{A thesis submitted for the degree of Master of Computer Science at École polytechnique fédérale de Lausanne}

  \date{June}{2016}

  \begin{center}
    \begin{tabular}{cc}
      \begin{tabular}{p{4.0cm}}
        \supervisor{Responsible}{Prof. Martin Odersky}{EPFL / LAMP}
      \end{tabular}&
      \begin{tabular}{p{4.0cm}}
        \supervisor{Supervisor}{Eugene Burmako}{EPFL / LAMP}
      \end{tabular}
    \end{tabular}
  \end{center}

% end title page




% Updated March 2006 (SP)


\newcommand{\logoepfl}[0]{
  \begin{center}
    \includegraphics[width=4cm]{logo_epfl_coul.eps}
  \end{center}
  \vspace{0.3cm}
  \hrule
}
\newcommand{\logolasec}[0]{
  \vspace{1cm}
  \hrule
  \begin{center}
    \includegraphics[width=4.5cm]{logo_lasec_coul.eps}
  \end{center}
}
\newcommand{\project}[1]{
  \begin{center}
    \large{#1}
  \end{center}
  \vspace{1cm}
}
\newcommand{\department}[1]{
  \begin{center}
    \large{#1}
  \end{center}
}
\newcommand{\supervisor}[3]{
  \begin{center}
    \begin{normalsize}{
        \bf #1}\\#2\\#3
    \end{normalsize}
  \end{center}
}
\renewcommand{\author}[1]{
  \begin{center}
    \Large{#1}
  \end{center}
  \vspace{0.5cm}
}
\renewcommand{\title}[1]{
  \vspace{3cm}
  \begin{center}
    \huge{#1}
  \end{center}
  \vspace{1.7cm}
}
\renewcommand{\date}[2]{
  \begin{center}
    \normalsize{#1 #2}
  \end{center}
  \vspace{0.5cm}
}


\thispagestyle{empty}


% begin title page
  \logoepfl

  \title{scalafmt: opinionated \\ code formatter for Scala}

  \author{Ólafur Páll Geirsson}
  \department{School of Computer and Communication Sciences}
  \project{A thesis submitted for the degree of Master of Computer Science at École polytechnique fédérale de Lausanne}

  \date{June}{2016}

  \begin{center}
    \begin{tabular}{cc}
      \begin{tabular}{p{4.0cm}}
        \supervisor{Responsible}{Prof. Martin Odersky}{EPFL / LAMP}
      \end{tabular}&
      \begin{tabular}{p{4.0cm}}
        \supervisor{Supervisor}{Eugene Burmako}{EPFL / LAMP}
      \end{tabular}
    \end{tabular}
  \end{center}

% end title page


\begin{abstract}
  Automatic code formatters bring a lot of benefits to software development.
  When done right, code formatters relieve the developer's attention from
  syntactic trivia and enforce a consistent coding style between developers.
  However, developing a good code formatter is still somewhat of a black art and yet,
  little research has been made towards the algorithms that power such code formatters.

  This thesis addresses the problem of developing a feature-rich code formatter for a custom programming language.
  % A critical feature of such a formatter is enforcing a column-width limit, which turns out to be a hard problem.
  Our contributions are twofold.
  First, we present a language agnostic framework that consists of core data structures, algorithms and tooling 
  that allow rapid development and testing of such a code formatter.
  Secondly, we use our model in a case study and provide a concrete implementation of a code formatter, \texttt{scalafmt}.
  We show that the framework makes


\end{abstract}

\begin{abstract}
  Automatic code formatters bring many benefits to software development, yet they can be tricky to implement.
  This thesis addresses the problem of developing a code formatter for the
  Scala programming language that captures many popular coding styles.
  Our work has been limited to formatting Scala code.
  Still, we have developed algorithms and tools,
  which we believe can be of interest to developers of code formatters for other programming languages.
\end{abstract}
\tableofcontents
% 
\section{Introduction} % (fold)
\label{sec:Introduction}
The main contributions presented in this thesis are the following:
\begin{itemize}
    \item A DSL for declaring feature-rich code formatting rules.
    \item An case study where the DSL is used to format Scala programs.
\end{itemize}
Throughout the paper we assume familiarity with the basics of the Scala
Programming Language~\autocite{odersky_scala_2004}.

\pagebreak
\section{Introduction} % (fold)
\label{sec:Introduction}
\lstset{style=scala}
The main motivation of this study is to bring scalafmt, a new Scala code formatter, to the Scala community.
The goal is to capture many popular coding styles so that a wide part of the Scala community can enjoy the benefits that come with automatic code formatting.

Without code formatters, software developers are responsible for manipulating all syntactic trivia in their programs.
What is syntactic trivia?
Consider the Scala code snippets in listings~\ref{lst:unformatted} and~\ref{lst:formatted1}.

\begin{minipage}{.45\textwidth}
\lstinputlisting[label={lst:unformatted}, caption=Unformatted code]{code/unformatted.scala}
\end{minipage}
\hfil
\begin{minipage}{.45\textwidth}
\lstinputlisting[label={lst:formatted1}, caption=Formatted code]{code/target/formatted1.scala}
\end{minipage}
% \hfil
Both snippets represent the same program.
The only difference lies in their syntactic trivia, that is where spaces and line breaks are used.
% The snippet in listing~\ref{lst:formatted1} is result of running \scalafmt{} on listing~\ref{lst:unformatted} with a constrait of fitting maximum 35 characters on each line.
Although the whitespace does not alter the execution of the program, listing~\ref{lst:formatted1} is arguably easier to read, understand and maintain for the software developer.
The promise of code formatters is to automatically convert any program that may contain style issues, such as in listing~\ref{lst:unformatted}, into a readable and consistent looking program, such as in listing~\ref{lst:formatted1}.
Automatic code formatting offers several benefits.
% Without a code formatters, the software developer is left to come up with the layout she prefers the most.
% With a code formatter, the choice is 
% In this example, good code formatting offers great value to the software developer.

Code formatting enables large-scale refactoring.
Google used ClangFormat\autocite{jasper_clang-format_????}, a code formatter, to migrate legacy C++ code to the modern C++11 standard\autocite{wright_large-scale_2013}.
ClangFormat was used to ensure that the refactored code adhered to Google's extensive C++ coding style\autocite{_google_????}.
Similar migrations can be expected in the near future for the Scala community once new dialects, such as Dotty\autocite{rompf_f_2015}, gain popularity.

Code formatting is valuable in collaborative coding environments.
The Scala.js project\autocite{_scala.js_????} has over 40 contributors and the Scala.js coding style\autocite{_scala.js_????-1} contains over 2.600 words.
Each contributor to the Scala.js project is expected to follow the coding style.
Each contributed patch is manually verified against the coding style by the project maintainers.
This adds a burden on both contributors and maintainers.
Several maintainers of popular Scala libraries have expressed this sentiment.
ENSIME\autocite{_ensime_????} is a popular Scala interaction mode for text editor.
Sam Halliday, a maintainer of ENSIME, says ``I don't have time to talk about formatting in code reviews. I want the machine to do it so I can focus on the design.''\autocite{halliday_i_2016-1}.
Akka\autocite{_akka_????} is another Scala library to build concurrent and distributed applications.
Viktor Klang, a maintainer of Akka, suggests a better alternative ``Code style should not be enforced by review, but by automate rewriting. Evolve the style using PRs against the rewriting config.''.\autocite{klang_code_2016}.
With code formatters, software develoeprs can direct their full focus on writing correct, maintainable and fast code.

\subsection{Whitespace style issues}
Mechanical style issues are source code issues that can be fixed automatically.
For example, consider listing~\ref{lst:mechanic}.
\lstinputlisting[label={lst:mechanic}, caption=Style issues]{code/mechanic.scala}
This code snippet leaves a lot left to be desired.
Mechanical issues include the redundant spaces around parentheses, unnecessary semicolon after the \texttt{println} and the inconsistent indentation in the body of the if statement.
Non-mechanical issues may include the fact the program prints an error message to the standard output and exits the process, instead of throwing an exception.

This thesis only addresses the whitespace style issues.
That is, to automatically fix the spaces and newline characters between the non-whitespace tokens in the original source code.
We will leave it to the software developer to decide whether the unnecessary semicolon should remain in the source file or if the program should throw an exception.
% Another example, the Spark style guide\autocite{spark-scala} follows the ``Rule of 30''\autocite{software_refactoring}.
% The rule of 30 dictates amongst several things that method bodies should not contain more than 30 lines of code.
% Although it is easy to statically detect a violation of that rule, it is non-trivial to automatically rewrite a method into smaller pieces of code.



% Code formatters have existed for a long time.
% In more recent years we've seen many new opimization based code formatters, such as clang-format, dartfmt and rfmt.
% Optimization based code formatters define a fit function for a given layout.
% A good layout has low cost while a bad layout has a high cost.
% The formatter runs an algorithm to choose the layout with the lowest code.
% This idea is not new, for example TeX uses a fit function to optimally break lines
% in text documents.

% Much work has been put into developing language agnostic code formatters or pretty printers.
% Researchers of such work correctly highlights that developing a custom code formatter for each language requires a lot of effort.
% It may be possible that we find a perfect way to develop language independent code formatters, see~\autocite{mps_article}.
% Still, today no such framework exists and people who write Scala code each day are still choosing to manually format their code.
% The goal of this thesis is not to address language agnostic code formatting.
% \subsection{Non-mechanical style issues}
% TODO(olafur)

% \subsection{Problem statement}
% This thesis addresses the problem of developing a code formatter, \scalafmt, for the Scala programming language.
% This formatter should:
% \begin{itemize}
%   \item capture Scala's most common idioms and popular coding styles, most importanly a line length limit.
%   \item be opinionated, it should be able to produce nice looking code even for the most egregiously formatted input.
%   \item be fast, <200ms for 99 percentile of source files, <2s for rest.
% \end{itemize}

\subsection{Contributions}
The main contribution presented in this thesis are the following:
\begin{itemize}
  \item scalafmt, a code formatter for the Scala programming language.
    At the time of this writing, scalafmt has been available for 3 months, it has been installed over 5.000 times and is already in use by several open source Scala libraries.
    For details on how to install and use scalafmt, refer to the scalafmt online documentation\autocite{geirsson_scalafmt_????}.
    % Discussion of scalafmt adoption is presented in section~\ref{sec:adoption}.

  \item algorithms and data structures to implement line wrapping under a maximum column-width limit.
    This work is presented in section~\ref{sec:algorithms}.
  \item tools to develop and test code formatters.
    This work is presented in section~\ref{sec:tooling}.
\end{itemize}
The scalafmt formatter itself may only be of direct interest to the Scala community.
However, much work in this thesis is not specifically tied to Scala and we hope can inspire the design of code formatters for other programming languages.

\section{Background}
This chapter explains the necessary background to understand Scala and code formatting.
More specifically, we motivate why Scala presents an interesting challenge for code formatters.
We go into details on Scala's rich syntax and popular idioms that introduced unique challenges to the design of scalafmt.
We follow up with a brief history on code formatters for both Scala as well as other programming languages.
We will see that although code formatters have a long history, the last 5 years have brought a lot of research in optimization based formatters, which scalafmt's design takes inspiration from.

\subsection{Scala the programming language}
Scala\autocite{odersky_scala_2004} is a general purpose programming language that was first released in 2004.
Scala combines features from object-oriented and functional programming paradigms, allowing maximum code reuse and extensibility.

Scala code can run on multiple platforms.
Most commonly, Scala programs compile to bytecode and run on the JVM.
With the releases of Scala.js\autocite{_scala.js_????-1}, JavaScript has recently become a popular target platform for Scala developers.
Even more recently, the announcement of Scala Native\autocite{_scala-native/scala-native_????} shows that LLVM and may become yet another viable platform for Scala developers.

Scala is a popular programming language.
The Scala Center estimates that more than half a million developers are using Scala\autocite{odersky_scala_2016}.
Large organizations such as Goldman Sachs, Twitter, IBM and Verizon run Scala code in production systems.
The 2015 Stack Overflow Developer Survey shows that Scala is the 6th most loved technology and 4th best paying technology to work with\autocite{_stack_????}.
The popularity of Apache Spark\autocite{_apache_????-1}, a cluster computing framework for large-scala data processing, has made Scala a language of choice for many developers and scientists working in big data and machine learning.

Scala is a programming language with rich syntax and many idioms.
The following chapters discuss in detail several prominent syntactic features and idioms of Scala.
Most importantly, we highlight coding patterns that encourage developers to write larger statements instead of many small statements.
In section~\ref{sec:algorithms}, we explain why large statements introduce a challenge to code formatting.

\subsubsection{Higher order functions}
Higher order functions (HOFs) are a common concept in functional programming languages and mathematics.
HOFs are functions that can take other functions as arguments and can return functions as return values.
Languages that provide a convenient syntax to manipulate HOFs are said to make functions first-class citizens.

Functions are first-class citizens in Scala.
Consider listing~\ref{lst:hof}.
\lstinputlisting[label={lst:hof}, float, caption=Higher order functions]{code/hof.scala}
The method \texttt{twice} takes an argument \texttt{f}, which is a function from an integer to an integer.
The method returns a new function that will apply \texttt{f} twice to an integer argument.
This small example takes advantage of several syntactic conveniences provided by Scala.
For example, in line 2 the argument \texttt{\_ + 3} creates a new \texttt{Function[Int, Int]} object.
The function call \texttt{f(x)} is in fact sugar for the method call \texttt{f.apply(x)} on a \texttt{Function[Int, Int]} instance.
Listing~\ref{lst:hof_nosugar} shows an equivalent program to listing~\ref{lst:hof} without syntactic sugar.
\lstinputlisting[label={lst:hof_nosugar}, float, caption=Higher order functions without syntactic sugar]{code/hof-nosugar.scala}
Observe that what was expressed as a single statement in line 1 of listing~\ref{lst:hof} is expressed with multiple statements in lines 1 and 2 of listing~\ref{lst:hof_nosugar}.

\subsubsection{Immutability}
Functional programming encourages stateless functions which operate on immutable data structures and objects.
An immutable object is an object that once initialized, cannot be modified.
Immutability offers several benefits to software development in areas including concurrency and testing.
Listing~\ref{lst:immutable} shows an example of manipulating an immutable list.
\lstinputlisting[label={lst:immutable}, float, caption=Manipulating immutable list]{code/immutable.scala}
Note that each \texttt{map} and \texttt{filter} operation creates a new copy of the list with the modified contents.
The original list remains unchanged.
Listing~\ref{lst:mutable} show the equivalent operation using a mutable list.
\lstinputlisting[label={lst:mutable}, float, caption=Manipulating mutable list]{code/mutable.scala}
Observe that what was listing~\ref{lst:immutable} is a single statement while listing~\ref{lst:mutable} is multiple statements.

\subsubsection{SBT build configuration}
SBT\autocite{_sbt_????} is an interactive build tool used by many Scala projects.
SBT configuration files are written in \texttt{*.sbt} or \texttt{*.scala} files using Scala.
Although SBT configuration files use plain Scala, they typically use coding patterns which are different from traditional Scala programs.
Listing~\ref{lst:sbt} is an example project definition in SBT.
\lstinputlisting[label={lst:sbt}, float, caption=SBT project definition]{code/sbt.scala}
Observe that the project is defined as a single statement and makes extensive use of symbolic infix operators.
Due to the nature of build configurations, argument lists to can becomes unwieldy long and a single project statement can span up to dozens or even hundreds of lines.

% \subsection{Metaprogramming with scala.meta}
% Scala.meta\autocite{_scala.meta_????} is a robust and portable metaprogramming toolkit for Scala.
% Scala 2.10 compiler introduced the first metaprogramming API for Scala via compile-time macros and runtime reflection.
% Library authors could take advantage of this new functionality to implement more advanced features in their libraries.
% However, the Scala compiler metaprogramming facilities exposed compiler internals to its users.

\subsection{Code formatting}
\subsubsection{History}
The science of displaying aesthetically pleasing text predates as early as 1956\autocite{harris_keyboard_1956}.
The first efforts involved printing natural language text.
As soon as higher-level programming languages evolved, programmers have researched how to appealingly display code.

Scowen\autocite{scowen_soapprogram_1971} developed SOAP in 1971, a code formatter for ALGOL 60.
The main motivation for SOAP was to make it ``easier for a programmer to examine and follow a program'' as well as to maintain a consistent coding style.
SOAP did provide a line length limit.
However, SOAP would fail execution if the provided line length turned out to be too small.
With the hardware at the time, SOAP could format 600 lines of code per minute.

Since then, each new 

\subsubsection{gofmt}
\subsubsection{rustfmt}
\subsubsection{dartfmt}
\subsubsection{clang-format}
\subsubsection{scalariform}
There already exists a widely used code formatter for Scala called Scalariform\autocite{scalariform}.
However, Scalariform lacks the ability to enforce a column width limit,
which we consider necessary to capture many popular Scala coding styles.
\section{Algorithms}\label{sec:algorithms}
\subsection{Data structures}
\subsubsection{FormatToken}
\subsubsection{Split}
\subsubsection{Policy}
\subsubsection{OptimalToken}
\subsubsection{State}
\subsubsection{Indent}
\subsection{BestFirstSearch}
\subsubsection{Router}
\begin{itemize}
  \item Router
  \item Search.
\end{itemize}
\subsection{Optimizations}
\subsubsection{dequeueOnNewStatements}
\subsubsection{recurseOnBlocks}
\subsubsection{escapeInPathologicalCases}
\subsubsection{escapeInPathologicalCases}
\subsubsection{pruneSlowStates}
\subsubsection{FormatWriter}
\begin{itemize}
  \item vertical alignment
  \item comment formatting
  \item stripMargin alignment
\end{itemize}
\section{Tooling}\label{sec:tooling}
\subsection{Heatmaps}
\subsection{Configuration}
\subsubsection{maxColumn}
\subsubsection{binPacking}
\subsubsection{vertical alignment}
\subsection{Unit tests}
\subsection{Property based tests}
\subsubsection{AST Integrity}
\subsubsection{Idempotency}
\subsection{Regressions tests}
\section{Evaluation}
\subsection{Micro benchmarks}
\subsection{Adoption}\label{sec:adoption}
% \section{Related work}
% \subsection{Combinator based}
% \begin{enumerate}
%   \item Houghes 1995
%   \item Wadler 1999
% \end{enumerate}
% \subsection{Optimization-oriented}
% \begin{enumerate}
%   \item clang-format Dijkstra's 2010
%   \item dartfmt Best-first search 2014
%   \item rfmt 2015
% \end{enumerate}
% \begin{enumerate}
%   \item Optimal line breaking
%   \item Oppen
% \end{enumerate}
\section{Discussion}
\subsection{Future work}
\subsection{Conclusion}
\printbibliography{}

\end{document}


